\documentclass[10pt,a4paper]{amsart}
\usepackage{amsmath,amssymb,amsthm,graphicx}
% Retained preamble for the Stacks Project chapter snapshot in this directory.
% The chapter's own preamble.tex is a copy of the shared Stacks Project preamble; the block
% below is sliced MECHANICALLY, line for line, from it. Nothing was added, removed or
% redefined; the wrapper class, the xr-hyper/hyperref wiring and the externaldocument list
% are omitted because the assembler supplies the document class.
%
% lines taken from preamble.tex: 3, 5, 9, 11, 12, 156, 157, 158, ... (51 lines total)
%
% package coverage restored from the chapter's own preamble (2026-10-13): xy, xy-2cell

\usepackage[all]{xy}
\xyoption{2cell}
\UseAllTwocells

\newenvironment{reference}{\comment}{\endcomment}
\newenvironment{slogan}{\comment}{\endcomment}
\newenvironment{history}{\comment}{\endcomment}
\theoremstyle{plain}
\newtheorem{theorem}[subsection]{Theorem}
\newtheorem{proposition}[subsection]{Proposition}
\newtheorem{lemma}[subsection]{Lemma}
\theoremstyle{definition}
\newtheorem{definition}[subsection]{Definition}
\newtheorem{example}[subsection]{Example}
\newtheorem{exercise}[subsection]{Exercise}
\newtheorem{situation}[subsection]{Situation}
\theoremstyle{remark}
\newtheorem{remark}[subsection]{Remark}
\newtheorem{remarks}[subsection]{Remarks}
\numberwithin{equation}{subsection}
\def\lim{\mathop{\mathrm{lim}}\nolimits}
\def\colim{\mathop{\mathrm{colim}}\nolimits}
\def\Spec{\mathop{\mathrm{Spec}}}
\def\Hom{\mathop{\mathrm{Hom}}\nolimits}
\def\Ext{\mathop{\mathrm{Ext}}\nolimits}
\def\SheafHom{\mathop{\mathcal{H}\!\mathit{om}}\nolimits}
\def\SheafExt{\mathop{\mathcal{E}\!\mathit{xt}}\nolimits}
\def\Sch{\mathit{Sch}}
\def\Mor{\mathop{\mathrm{Mor}}\nolimits}
\def\Ob{\mathop{\mathrm{Ob}}\nolimits}
\def\Sh{\mathop{\mathit{Sh}}\nolimits}
\def\NL{\mathop{N\!L}\nolimits}
\def\CH{\mathop{\mathrm{CH}}\nolimits}
\def\proetale{{pro\text{-}\acute{e}tale}}
\def\etale{{\acute{e}tale}}
\def\QCoh{\mathit{QCoh}}
\def\Ker{\mathop{\mathrm{Ker}}}
\def\Im{\mathop{\mathrm{Im}}}
\def\Coker{\mathop{\mathrm{Coker}}}
\def\Coim{\mathop{\mathrm{Coim}}}
\DeclareMathSymbol{\boxtimes}{\mathbin}{AMSa}{"02}
\def\QCohstack{\mathcal{QC}\!\mathit{oh}}
\def\Cohstack{\mathcal{C}\!\mathit{oh}}
\def\Spacesstack{\mathcal{S}\!\mathit{paces}}
\def\Quotfunctor{\mathrm{Quot}}
\def\Hilbfunctor{\mathrm{Hilb}}
\def\Curvesstack{\mathcal{C}\!\mathit{urves}}
\def\Polarizedstack{\mathcal{P}\!\mathit{olarized}}
\def\Complexesstack{\mathcal{C}\!\mathit{omplexes}}
\def\Pic{\mathop{\mathrm{Pic}}\nolimits}
\def\Picardstack{\mathcal{P}\!\mathit{ic}}
\def\Picardfunctor{\mathrm{Pic}}
\def\Deformationcategory{\mathcal{D}\!\mathit{ef}}
% lane-s11: discard bibliographic \begin{reference} blocks entirely (incl. any \cite inside).
\newsavebox{\limrefbox}
\ifcsname reference\endcsname\expandafter\let\csname reference\endcsname\undefined\fi\ifcsname endreference\endcsname\expandafter\let\csname endreference\endcsname\undefined\fi
\def\citation#1{\relax}
\newenvironment{reference}{\setbox\limrefbox=\vbox\bgroup\let\cite\relax}{\egroup\ignorespaces}
\usepackage{cleveref}
\crefname{originalproposition}{original proposition}{original propositions}
\Crefname{originalproposition}{Original proposition}{Original propositions}
\makeatletter
\newlabel{proposition-characterize-locally-finite-presentation}{{01ZC}{source}{}{sourcepointer-1}{}}
\newlabel{proposition-characterize-locally-finite-presentation@cref}{{[originalproposition][0][]01ZC}{[1][0][]source}{}{}{}}
\makeatother
\crefname{originalreference}{original reference}{original references}
\Crefname{originalreference}{Original reference}{Original references}
\makeatletter
\newlabel{algebra-lemma-characterize-finite-presentation}{{00QO}{source}{}{sourcepointer-2}{}}
\newlabel{algebra-lemma-characterize-finite-presentation@cref}{{[originalreference][0][]00QO}{[1][0][]source}{}{}{}}
\makeatother
\crefname{originallemma}{original lemma}{original lemmas}
\Crefname{originallemma}{Original lemma}{Original lemmas}
\makeatletter
\newlabel{lemma-limit-closed-nonempty}{{01Z3}{source}{}{sourcepointer-3}{}}
\newlabel{lemma-limit-closed-nonempty@cref}{{[originallemma][0][]01Z3}{[1][0][]source}{}{}{}}
\makeatother
\title{Surjectivity of the Mor map suffices for the colimit identity}
\author{The Stacks Project Authors}
\date{stacks-project:limits.tex:local-snapshot-20261012; GNU FDL 1.2 or later}
\begin{document}
\maketitle
\noindent\textit{Editorial note.} No omissions: the argument is fully repeated in the counted span rather than citing the characterization proposition.\par\medskip
\noindent\textit{Allowed standard prerequisites (editorial).}\begin{itemize}
\item Limit closed nonempty. Limits, Lemma lemma-limit-closed-nonempty
\end{itemize}\medskip
\ifdefined\hypertarget\hypertarget{sourcepointer-1}{}\fi
\noindent\textit{Source reference.} Original proposition 01ZC: sources/stacks-chow/tags.txt, source lines 2587--2587. Cross-chapter standard prerequisite cited by the author inside the counted span; it is used as a pointer and its own proof is not part of this item.\par\medskip
\ifdefined\hypertarget\hypertarget{sourcepointer-2}{}\fi
\noindent\textit{Source reference.} Original reference 00QO: sources/stacks-chow/tags.txt, source lines 974--974. Cross-chapter standard prerequisite cited by the author inside the counted span; it is used as a pointer and its own proof is not part of this item.\par\medskip
\ifdefined\hypertarget\hypertarget{sourcepointer-3}{}\fi
\noindent\textit{Source reference.} Original lemma 01Z3: sources/stacks-chow/tags.txt, source lines 2578--2578. Cross-chapter standard prerequisite cited by the author inside the counted span; it is used as a pointer and its own proof is not part of this item.\par\medskip
\begin{lemma}
\label{lemma-surjection-is-enough}
Let $f : X \to S$ be a morphism of schemes. If for every directed limit
$T = \lim_{i \in I} T_i$ of affine schemes over $S$ the map
$$
\colim \Mor_S(T_i, X) \longrightarrow \Mor_S(T, X)
$$
is surjective, then $f$ is locally of finite presentation.
In other words, in
Proposition \ref{proposition-characterize-locally-finite-presentation}
parts (2) and (3) it suffices to check surjectivity of the map.
\end{lemma}

\begin{proof}
The proof is exactly the same as the proof of the implication
``(2) implies (1)'' in
Proposition \ref{proposition-characterize-locally-finite-presentation}.
Choose any affine opens $U \subset X$ and $V \subset S$ such that
$f(U) \subset V$. We have to show that
$\mathcal{O}_S(V) \to \mathcal{O}_X(U)$ is of finite presentation.
Let $(A_i, \varphi_{ii'})$ be a directed system of
$\mathcal{O}_S(V)$-algebras. Set $A = \colim_i A_i$.
According to
Algebra, Lemma \ref{algebra-lemma-characterize-finite-presentation}
it suffices to show that
$$
\colim_i \Hom_{\mathcal{O}_S(V)}(\mathcal{O}_X(U), A_i) \to
\Hom_{\mathcal{O}_S(V)}(\mathcal{O}_X(U), A)
$$
is surjective. Consider the schemes $T_i = \Spec(A_i)$. They
form an inverse system of $V$-schemes over $I$
with transition morphisms $f_{ii'} : T_i \to T_{i'}$
induced by the $\mathcal{O}_S(V)$-algebra maps $\varphi_{i'i}$.
Set $T := \Spec(A) = \lim_i T_i$.
The formula above becomes in terms of morphism sets of schemes
$$
\colim_i \Mor_V(T_i, U) \to \Mor_V(\lim_i T_i, U)
$$
We first observe that
$\Mor_V(T_i, U) = \Mor_S(T_i, U)$
and
$\Mor_V(T, U) = \Mor_S(T, U)$.
Hence we have to show that
$$
\colim_i \Mor_S(T_i, U) \to
\Mor_S(\lim_i T_i, U)
$$
is surjective and we are given that
$$
\colim_i \Mor_S(T_i, X) \to
\Mor_S(\lim_i T_i, X)
$$
is surjective.
Hence it suffices to prove that given a morphism $g_i : T_i \to X$ over $S$
such that the composition $T \to T_i \to X$ ends up in $U$ there exists some
$i' \geq i$ such that the composition $g_{i'} : T_{i'} \to T_i \to X$ ends up
in $U$. Denote $Z_{i'} = g_{i'}^{-1}(X \setminus U)$.
Assume each $Z_{i'}$ is nonempty
to get a contradiction. By Lemma \ref{lemma-limit-closed-nonempty}
there exists a point $t$ of $T$ which is mapped into $Z_{i'}$ for all
$i' \geq i$. Such a point is not mapped into $U$. A contradiction.
\end{proof}

\end{document}
